% Problem record op_c45752aeb9a106bc. This TeX file is derived from
% database/problems_json/op_c45752aeb9a106bc.json by scripts/sync-tex.mjs; edit the JSON record
% and rerun the script rather than editing this file.
\section{Channels superactivated by antidegradable helpers}
\label{sec:op_c45752aeb9a106bc}

\paragraph{Problem.}
What necessary and sufficient structural conditions characterize the finite-dimensional channels of zero quantum capacity, and separately those of zero private capacity, that acquire a positive capacity of the same kind when used in parallel with some antidegradable channel?
Let $Q(\mathcal N)$ and $P(\mathcal N)$ denote the fully regularized unassisted quantum capacity and private classical capacity of a finite-dimensional memoryless quantum channel $\mathcal N$. Rates are measured in qubits and private bits per channel use, respectively; the eavesdropper receives the full complementary-channel output. For an isometric dilation $V:A\to B\otimes E$, define the channel and its complementary channel by
\begin{equation}
 \mathcal N(\rho):=\operatorname{Tr}_E(V\rho V^\dagger),
 \qquad
 \mathcal N^c(\rho):=\operatorname{Tr}_B(V\rho V^\dagger).
 \label{eq:c457-dilation}
\end{equation}
A channel $\mathcal A$ is antidegradable when $\mathcal A=\mathcal D\circ\mathcal A^c$ for some completely positive trace-preserving map $\mathcal D$, with $\mathcal A^c$ defined as in Eq.~\eqref{eq:c457-dilation}. Let $\mathfrak A$ denote the set of all finite-dimensional antidegradable channels, with no fixed dimension bound. For $R\in\{Q,P\}$, define the activation set
\begin{equation}
 \mathsf{Act}_R:=\bigl\{\mathcal N:\ R(\mathcal N)=0\ \text{and}\ R(\mathcal N\otimes\mathcal A)>0\ \text{for some}\ \mathcal A\in\mathfrak A\bigr\}.
 \label{eq:c457-activation-sets}
\end{equation}
The product $\mathcal N\otimes\mathcal A$ uses independent channel environments and allows arbitrary joint encoding and decoding across uses, but no public communication assistance. The problem asks, for $R=Q$ and for $R=P$, for a necessary and sufficient structural condition on $\mathcal N$ for membership in the set defined in Eq.~\eqref{eq:c457-activation-sets}.

\subsection*{Status}

Unsolved

\subsection*{Source}

The superactivation paper of Smith and Yard closes by asking, among other questions, ``Are there other pairs of zero-capacity channels displaying this effect?'', ``Does the private capacity also display superactivation?'', and ``Can all Horodecki channels be superactivated, or just those with positive private capacity?'' \sourcecite{ref:c457-smith-yard}{SY08}. In Section~4 of the arXiv version, Hirche and Leditzky note that private-capacity superactivation requires a factor that is not antidegradable, and suggest channels that are PPT, have PPT complements, and are not antidegradable as candidates for private-capacity superactivation together with antidegradable channels \sourcecite{ref:c457-hirche-leditzky}{HL23}. The outlook of Zhu and Wang's September 2026 preprint asks which zero-private-capacity channels admit the support asymmetry used in its activation protocol \sourcecite{ref:c457-zhu-wang-private}{ZW26b}. The characterization question for antidegradable helpers, posed separately for quantum and private capacity, is a synthesis of these structural questions. No cited source conjectures that every channel outside a known zero-capacity class is activatable.

\subsection*{Progress}

\begin{itemize}
  \item Smith and Yard proved, in the theorem of their Supporting Online Material, that every channel $\mathcal C$ and every finite input ensemble $\{p_u,\rho_u\}$ satisfy
\begin{equation}
 Q(\mathcal C\otimes\mathcal E_{d,1/2})\geq\frac{\Delta}{2},
 \qquad
 \Delta:=I(U:B)-I(U:E),
 \qquad
 \mathcal E_{d,1/2}(\rho):=\frac12\rho\oplus\frac12\operatorname{Tr}(\rho)|e\rangle\langle e|.
 \label{eq:c457-smith-yard}
\end{equation}
Here $U$ is the classical ensemble label, $I$ is quantum mutual information evaluated on the output $B$ of $\mathcal C$ and the output $E$ of $\mathcal C^c$, the erasure channel acts on $\mathbb C^d$ with $d$ the sum of the ranks of the states $\rho_u$, and $|e\rangle$ is an orthogonal erasure flag. The 50\% erasure channel $\mathcal E_{d,1/2}$ is antidegradable. For a four-dimensional channel that is PPT, meaning that transposition of its output composed with the channel is completely positive, and an ensemble with $\Delta>0.02$, Eq.~\eqref{eq:c457-smith-yard} gives a quantum capacity above $0.01$ with $d=4$, so $\mathsf{Act}_Q$ is nonempty \sourcecite{ref:c457-smith-yard}{SY08}. More generally, if $Q(\mathcal N)=0<P(\mathcal N)$, some tensor power $\mathcal N^{\otimes k}$ has an ensemble with $\Delta>0$. Applying Eq.~\eqref{eq:c457-smith-yard} to $\mathcal C=\mathcal N^{\otimes k}$ and using $kQ(\mathcal N\otimes\mathcal E_{d,1/2})\geq Q(\mathcal N^{\otimes k}\otimes\mathcal E_{d,1/2})$ shows that every such channel lies in $\mathsf{Act}_Q$.

  \item Two antidegradable channels cannot superactivate either capacity. If $\mathcal A_j=\mathcal D_j\circ\mathcal A_j^c$ for $j=1,2$, then
\begin{equation}
 \mathcal A_1\otimes\mathcal A_2=(\mathcal D_1\otimes\mathcal D_2)\circ(\mathcal A_1^c\otimes\mathcal A_2^c),
 \qquad
 Q(\mathcal A_1\otimes\mathcal A_2)=P(\mathcal A_1\otimes\mathcal A_2)=0,
 \label{eq:c457-antidegradable-products}
\end{equation}
because $\mathcal A_1^c\otimes\mathcal A_2^c$ is a complementary channel of the product. This observation, used in the proof of Corollary~4.1 of Hirche and Leditzky, holds for arbitrary block codes, not merely product inputs. It shows that $\mathfrak A$ is disjoint from both activation sets \sourcecite{ref:c457-hirche-leditzky}{HL23}. Vanishing single-use coherent or private information alone is not the corresponding zero-capacity condition.

  \item Zhu and Wang's July 2026 preprint proves, in its Theorem~1.1, with proofs in Theorem~4.3 and Section~5, that the qutrit channel
\begin{equation}
 \Lambda(X):=\frac12X+\frac14\bigl(\operatorname{Tr}(X)I_3-X^{T}\bigr)
 \qquad\text{satisfies}\qquad
 P(\Lambda)=Q(\Lambda)=0,
 \label{eq:c457-qutrit-channel}
\end{equation}
where $I_3$ is the qutrit identity and $T$ is matrix transposition, although $\Lambda$ is neither antidegradable nor PPT. The proof shows that $\Lambda^c$ dominates $\Lambda$ in the complete less-noisy order: for every finite-dimensional reference system $R$ and all states $\rho,\sigma$ on $RA$ with $\operatorname{supp}\rho\subseteq\operatorname{supp}\sigma$,
\begin{equation}
 D\bigl((\operatorname{id}_R\otimes\Lambda^c)(\rho)\,\big\|\,(\operatorname{id}_R\otimes\Lambda^c)(\sigma)\bigr)
 \geq
 D\bigl((\operatorname{id}_R\otimes\Lambda)(\rho)\,\big\|\,(\operatorname{id}_R\otimes\Lambda)(\sigma)\bigr),
 \label{eq:c457-complete-less-noisy}
\end{equation}
where $D$ is the Umegaki relative entropy \sourcecite{ref:c457-zhu-wang-qutrit}{ZW26a}.

  \item The order in Eq.~\eqref{eq:c457-complete-less-noisy} tensorizes, by Lemma~3.1 of Hirche, Rouz\'e, and Stilck Fran\c{c}a. By their Proposition~3.2, a channel dominates in this order every channel obtained from it by post-processing, so $\mathcal A^c$ dominates $\mathcal A$ for every $\mathcal A\in\mathfrak A$ \sourcecite{ref:c457-hirche-rouze-franca}{HRS22}. Hence $\Lambda^c\otimes\mathcal A^c$, a complementary channel of $\Lambda\otimes\mathcal A$, dominates $\Lambda\otimes\mathcal A$, and Corollary~3.1 of the July preprint yields
\begin{equation}
 P(\Lambda\otimes\mathcal A)=Q(\Lambda\otimes\mathcal A)=0
 \qquad\text{for every }\mathcal A\in\mathfrak A.
 \label{eq:c457-qutrit-obstruction}
\end{equation}
The discussion of Zhu and Wang's September 2026 preprint states this exclusion of antidegradable helpers \sourcecite{ref:c457-zhu-wang-qutrit}{ZW26a}, \sourcecite{ref:c457-zhu-wang-private}{ZW26b}. The same argument excludes from both activation sets every channel that is dominated by its complement in the complete less-noisy order. Consequently, lying outside the antidegradable and PPT classes is not sufficient for membership in either set.

  \item Zhu and Wang's September 2026 preprint proves, in its Theorem~1, private-capacity superactivation with an antidegradable helper. Its channel $\mathcal N_0$, with four-dimensional input and output and defined by the Kraus form in its Eq.~(2.2), satisfies
\begin{equation}
 P(\mathcal N_0)=P(\mathcal E_{2,1/2})=0,
 \qquad
 P(\mathcal N_0\otimes\mathcal E_{2,1/2})\geq\frac{3\ln2}{10927}>1.903\times10^{-4}
 \label{eq:c457-private-superactivation}
\end{equation}
private bits per product use, where $\mathcal E_{2,1/2}$ is the qubit case of the erasure channel in Eq.~\eqref{eq:c457-smith-yard}. Thus $\mathsf{Act}_P$ is nonempty, subject to the qualification that this is a recent preprint result; its capacity inequalities are reported to be formalized in Lean~4. The zero private capacity of $\mathcal N_0$ follows from a completely positive trace-preserving map $\mathcal D$ with $\mathcal D\circ\mathcal N_0^c=T\circ\mathcal N_0$, so the environment has at least the receiver's Holevo information at every blocklength. Because $\mathcal N_0$ is activated, the argument leading to Eq.~\eqref{eq:c457-qutrit-obstruction} shows that $\mathcal N_0^c$ does not dominate $\mathcal N_0$ in the complete less-noisy order. The rate bound compares the receiver's information from a fixed binary measurement with the full environment's Holevo information, for one product use and a two-letter ensemble whose states have ranks two and three. Proposition~3 of the preprint shows that decoders whose effects remain positive under partial transposition of the output of $\mathcal N_0^{\otimes n}$ achieve no positive private rate for such products. The preprint's outlook asks which zero-private-capacity channels admit the support asymmetry used in this construction \sourcecite{ref:c457-zhu-wang-private}{ZW26b}.

  \item Combining Eqs.~\eqref{eq:c457-smith-yard} and \eqref{eq:c457-private-superactivation} shows that zero private capacity of both factors does not preclude quantum-capacity superactivation. By data processing, the two-letter ensemble of the previous item has private-information advantage $\Delta\geq3\ln2/10927$ for $\mathcal C:=\mathcal N_0\otimes\mathcal E_{2,1/2}$, and the ranks of its states sum to $d=5$. For the antidegradable channel $\mathcal A:=\mathcal E_{2,1/2}\otimes\mathcal E_{5,1/2}$, Eq.~\eqref{eq:c457-smith-yard} gives
\begin{equation}
 P(\mathcal N_0)=P(\mathcal A)=0,
 \qquad
 Q(\mathcal N_0\otimes\mathcal A)\geq\frac{3\ln2}{21854}>0,
 \label{eq:c457-quantum-activation-zero-private}
\end{equation}
so $\mathcal N_0$ belongs to both $\mathsf{Act}_Q$ and $\mathsf{Act}_P$. Since $T\circ\mathcal N_0=\mathcal D\circ\mathcal N_0^c$ is completely positive, $\mathcal N_0$ is itself a PPT channel; hence Eq.~\eqref{eq:c457-quantum-activation-zero-private} shows that positive private capacity is not necessary for a PPT channel to belong to $\mathsf{Act}_Q$, which bears on the question of Smith and Yard about Horodecki channels quoted in Source. This corollary is deduced here from the two cited results and is conditional on the September preprint. It does not assert that $\mathcal N_0\otimes\mathcal E_{2,1/2}$ itself has positive quantum capacity \sourcecite{ref:c457-smith-yard}{SY08}, \sourcecite{ref:c457-zhu-wang-private}{ZW26b}.
\end{itemize}

\subsection*{References}

\begin{enumerate}[label={},leftmargin=4.8em,labelsep=0.5em]
  \footnotesize
  \item[\textup{[SY08]}]\label{ref:c457-smith-yard}
  G. Smith and J. Yard, ``Quantum Communication with Zero-Capacity Channels,'' \emph{Science} \textbf{321}, 1812--1815 (2008). \href{https://doi.org/10.1126/science.1162242}{doi:10.1126/science.1162242}; \href{https://arxiv.org/abs/0807.4935}{arXiv:0807.4935}.

  \item[\textup{[HL23]}]\label{ref:c457-hirche-leditzky}
  C. Hirche and F. Leditzky, ``Bounding Quantum Capacities via Partial Orders and Complementarity,'' \emph{IEEE Transactions on Information Theory} \textbf{69}, 283--297 (2023). \href{https://doi.org/10.1109/TIT.2022.3199578}{doi:10.1109/TIT.2022.3199578}; \href{https://arxiv.org/abs/2202.11688}{arXiv:2202.11688}.

  \item[\textup{[HRS22]}]\label{ref:c457-hirche-rouze-franca}
  C. Hirche, C. Rouz\'e, and D. Stilck Fran\c{c}a, ``On Contraction Coefficients, Partial Orders and Approximation of Capacities for Quantum Channels,'' \emph{Quantum} \textbf{6}, 862 (2022). \href{https://doi.org/10.22331/q-2022-11-28-862}{doi:10.22331/q-2022-11-28-862}; \href{https://arxiv.org/abs/2011.05949}{arXiv:2011.05949}.

  \item[\textup{[ZW26a]}]\label{ref:c457-zhu-wang-qutrit}
  C. Zhu and X. Wang, ``Quantum Incapacity beyond No-Cloning and PPT Mechanisms,'' arXiv preprint, version 2 (2026). \href{https://arxiv.org/abs/2607.24693}{arXiv:2607.24693}.

  \item[\textup{[ZW26b]}]\label{ref:c457-zhu-wang-private}
  C. Zhu and X. Wang, ``Private Communication via Zero-Private-Capacity Quantum Channels,'' arXiv preprint, version 1 (2026). \href{https://arxiv.org/abs/2609.10520}{arXiv:2609.10520}.
\end{enumerate}

\subsection*{Comment}

The literature checked through 15 September 2026 gives members of both activation sets and structural obstructions, but no necessary and sufficient characterization of either set. Known members of $\mathsf{Act}_Q$ include every channel with $Q=0<P$; conditionally on the September 2026 preprint, the PPT channel $\mathcal N_0$ of zero private capacity belongs to both sets. Known non-members of both sets include all antidegradable channels and all channels dominated by their complements in the complete less-noisy order, such as $\Lambda$. The cited results supply no criterion deciding membership for a general channel of zero private capacity that is not dominated by its complement in that order. The nonexistence of activation between two antidegradable channels is settled and is not an open subproblem.
By Corollary~4.1 of Hirche and Leditzky, private-capacity superactivation implies that degradable channels form a proper subset of the regularized less-noisy channels \sourcecite{ref:c457-hirche-leditzky}{HL23}. That separation, the question of \href{https://qiqc-op.com/problem/op_fd75613c5bab4164/}{strict inclusion of degradable channels in the less-noisy class}, was already settled through the complement of $\Lambda$; conditionally on the September 2026 preprint, the channel $\mathcal N_0^c$ gives a further example.
The two-way-assisted secret-key analogue is the record on
\href{https://qiqc-op.com/problem/op_59e9ede6e282a917/}{superactivation of two-way secret-key capacity}.

\subsection*{Field}

Quantum Communication; Quantum Cryptography

\subsection*{Topic}

Superactivation; Quantum capacity; Private capacity; Channel degradability

\subsection*{ID}

\texttt{op\_c45752aeb9a106bc}
